Współczesna analiza numeryczna oraz uczenie maszynowe coraz częściej mierzą się z problemem aproksymacji funkcji wielu zmiennych $f: \mathbb{R}^d \to \mathbb{R}$ dla bardzo dużych wymiarów $d$. Klasyczne podejście do tego zagadnienia napotyka jednak fundamentalną barierę. Jak wykazali Novak i Woźniakowski \cite{novak2008tractability}, bez przyjęcia dodatkowych założeń o strukturze funkcji liczba próbek wymagana do osiągnięcia zadanego błędu aproksymacji rośnie wykładniczo wraz z wymiarem $d$. Zjawisko to, znane jako klątwa wymiarowości (ang. \textit{curse of dimensionality}), sprawia, że znaczna część standardowych algorytmów staje się w praktyce nieefektywna obliczeniowo (ang. \textit{intractable}).

Konstrukcja efektywnych algorytmów wymaga zatem przyjęcia modelu, w którym funkcja istotnie zależy jedynie od niewielkiej liczby parametrów. Przykładem takiego podejścia jest klasa tzw. funkcji grzbietowych (ang. \textit{ridge functions}), wprowadzona pierwotnie w kontekście tomografii komputerowej przez Logana i Sheppa \cite{logan1975optimal}.

Niniejsza praca magisterska jest poświęcona analizie uogólnionego modelu funkcji postaci
$$f(x) = g(Ax),$$
gdzie $x \in \mathbb{R}^d$, macierz $A$ ma wymiar $k \times d$ (przy założeniu, że $k \ll d$), a $g$ jest pewną funkcją nieliniową. Model ten stanowi uogólnienie podejścia badanego m.in. przez Cohena i in. \cite{cohen2011capturing} w kontekście pojedynczych wektorów stochastycznych.

Głównym celem pracy jest analiza teoretyczna i numeryczna algorytmów uczenia wyżej wymienionych funkcji na podstawie niewielkiej liczby losowych próbek. Podstawę merytoryczną rozważań stanowi artykuł Fornasiera, Schnassa i Vybirala pt. \textit{Learning Functions of Few Arbitrary Linear Parameters in High Dimensions} \cite{fornasier2012learning}. Została w nim zaproponowana metoda pozwalająca na odtworzenie macierzy parametrów $A$ z prawdopodobieństwem bliskim jedności, przy użyciu liczby próbek zależnej wielomianowo od wymiaru $d$.

Analiza mechanizmu działania tych algorytmów wymaga zastosowania zaawansowanego aparatu matematycznego. Kluczową rolę odgrywa tu teoria próbkowania skompresowanego (ang. \textit{compressed sensing}), a w szczególności metody stabilnej rekonstrukcji sygnałów rzadkich opracowane przez Candèsa, Romberga i Tao \cite{candes2006stable}. W cytowanej pracy wykazano, że gradienty badanej funkcji wykazują strukturę rzadką. Pozwala to na ich odzyskanie przy użyciu macierzy losowych spełniających własność ograniczonej izometrii (ang. \textit{Restricted Isometry Property}, RIP). Dowód tej własności dla macierzy Bernoulliego podali Baraniuk i in. \cite{baraniuk2008simple}. Istotnym elementem analizy jest również zjawisko koncentracji miary na sferze wysokowymiarowej, szeroko opisane przez Ledoux \cite{ledoux2001concentration}. Gwarantuje ono, że losowo wybrane punkty z bardzo dużym prawdopodobieństwem dostarczają użytecznych informacji o funkcji.

Praca została podzielona na trzy rozdziały:

W rozdziale pierwszym wprowadzono niezbędne definicje i twierdzenia. Omówiono w nim geometrię wysokich wymiarów oraz podstawy teorii próbkowania skompresowanego. Przedstawiono również kluczowe nierówności probabilistyczne, w tym macierzowe nierówności Chernoffa w ujęciu Troppa \cite{tropp2012user}, które są niezbędne do szacowania wartości osobliwych sum macierzy losowych w przypadku wielowymiarowym (dla $k > 1$).

Rozdział drugi zawiera szczegółowe omówienie analizy teoretycznej zaprezentowanej w artykule \cite{fornasier2012learning}. Opisano w nim algorytmy dla przypadku $k=1$ oraz ich uogólnienie dla $k>1$ oparte na rozkładzie SVD. Do analizy stabilności podprzestrzeni wykorzystano klasyczne twierdzenie Wedina \cite{wedin1972perturbation}, pozwalające oszacować błąd rekonstrukcji podprzestrzeni rozpiętej przez wiersze macierzy $A$.

Rozdział trzeci koncentruje się na aspektach praktycznych i implementacyjnych. Przedstawiono w nim wyniki eksperymentów numerycznych badających skuteczność algorytmów w obecności szumu oraz ich stabilność numeryczną. Implementację operacji macierzowych, takich jak rozkład SVD, oparto na standardowych metodach numerycznych opisanych w \cite{golub1996matrix}. Rozdział ten stanowi weryfikację teoretycznych oszacowań złożoności i zbieżności w rzeczywistych warunkach obliczeniowych.